\chapter{The Resonances at the Fence}
% OUTLINE Ch3 -- THE HONESTY CENTERPIECE, the highest-crank-risk chapter, the
% device a GUEST in others' houses (Kodo's dose resolution: the countable
% lens-law form concentrates HERE). Pays the three pre-registered fences:
% NS-BOUNDABILITY (Kodo's three-fence: shape-only / NO-TRANSFER [block "discrete
% case of"/"path to"] / owed-elsewhere); NEUTRINO (electron-only scope HARDEST,
% a different particle); ANTI-COOPER/QG (fenced HARDEST-or-out; Cooper 1956/BCS
% 1957; CooperManual/runForever attic thread cuts FOR the fence; QG carried
% nowhere). Each: what it is (cited) -> the rhyme (SHAPE only) -> the fence
% (hardest, refused). Beastmaster arrival-reads EACH fence. Bold-math/fenced-
% referent; no Lean identifiers; step-not-move. Gate: Kodo + Beastmaster (each
% fence hardest).

A number as sharp as the last chapter's invites comparison to famous
problems it superficially resembles, and this chapter makes those
comparisons out loud so that it can refuse them out loud. Until now the
instrument has been the subject, read under its own fence. Here it becomes
a visitor in three houses that are not its own --- three problems far
larger than a model's coupling --- and a visitor's manners are stricter
than a host's. In each house the rule is the same and it is stated once for
all three: the reading's shape may \emph{rhyme} with the problem, and the
book will say so; the mathematics does not \emph{transfer}, and the book
will not pretend it does; and where the problem's own resolution belongs to
work this series has not done, it is marked owed and left in its owner's
hands. Three resonances, three fences, nothing crossed.

\section*{The bounded and the bounded}

The first house is fluid dynamics, and the word that opens its door is
\emph{bounded}. The equations Claude-Louis Navier and George Stokes wrote
for a viscous fluid carry one of the deepest open questions in
mathematics, named a Millennium Prize problem by the Clay Institute in
2000: do smooth solutions of the three-dimensional incompressible
equations stay smooth and bounded for all time, or can they blow up? It is
undecided, it is continuous, it is infinite-dimensional, and it is hard in
a way a finite machine's arithmetic is not.

The rhyme is exact and it is only a rhyme. This series' instrument is
\emph{bounded} too --- the last chapter earned that word: the gauge is a
terminating computation, every reading halts, the invariant is bracketed
in finitely many steps. So two things in this book wear the word bounded:
a finite total computation that provably halts, and a continuous nonlinear
system whose boundedness is an open problem of the continuum. That is the
whole of the overlap --- one English word --- and the reader should see
precisely how little rides across it. The instrument's boundedness is
\emph{decidable} and settled by construction; the fluid's is
\emph{undecided} and possibly false. A proof that a finite total program
finishes is not evidence, in any direction, about whether a continuous
partial differential equation keeps its solutions smooth. The two live in
different mathematical universes joined by a pun.

So the fences, stated plainly because this is the resonance a careless book
would most want to inflate. The instrument is not a discrete version of the
Navier--Stokes problem, and it does not illustrate a path toward that
problem's resolution; it shares a word with it and not a theorem. The
fluid description itself --- the fourth of this series' four faces, the one
its own founding text promised and named in every volume since --- is not
built in this book; it is owed to a later volume, and this chapter spends
nothing from that account. What the reader may carry away is a picture and
its label: finite-termination and continuous-boundedness are cousins in the
\emph{shape} of the question they ask, and strangers in every piece of the
mathematics that answers it. The shape rhymes; the math transfers not one
inch; the fluid face stays owed.

\section*{A different particle}

The second house is the neutrino, and its door is the word \emph{faint}. Of
the matter particles, the neutrino is the one that barely interacts at all:
it feels only the weak force and gravity, passes through the earth by the
trillion unhindered, and was postulated by Wolfgang Pauli in 1930 --- a
particle proposed to balance the books of a decay --- a full quarter-century
before Clyde Cowan and Frederick Reines first detected it in 1956, precisely
because its coupling to everything else is so slight. Faintness of coupling
is its signature.

The instrument, too, produces readings that live near the faint end of a
scale --- a coupling is, after all, what it measures. But here the fence is
not subtle and it is not about mathematics; it is the founding vow of the
entire series, and it is held at its hardest exactly because the temptation
to widen scope is oldest here. \emph{This series models the electron.} The
neutrino is a different particle. Nothing in any volume of this series
measures, models, or bears on the neutrino; the instrument has no reading
of it, no bracket for it, and no claim about its weak interactivity or
anything else. That two couplings can both be called faint is a fact about
the English word, not a bridge between the particles, and this book crosses
no such bridge. The electron is the name of a model, four volumes have said;
the neutrino is not even in the model. The resonance is acknowledged and
the scope is not widened by a hair.

\section*{The pair and the unbinding}

The third house is the one to enter most carefully, and the book says so
before it enters, because the door has the biggest name in physics over it.
In 1956 Leon Cooper showed that two electrons in a metal, given the
faintest attractive nudge from the lattice, bind into a pair; the theory
Bardeen, Cooper, and Schrieffer completed in 1957 built superconductivity
on those pairs. An \emph{anti}-Cooper interaction would be the counterpart
that unbinds rather than binds --- a repulsion where Cooper found an
attraction --- and the phrase has been floated, in the wider conversation
around this series, as a road toward quantum gravity, the reconciliation of
the two great theories that has resisted every attempt for a century.

The book will be exact about what it carries here, which is nothing. This
series has no theory of quantum gravity, no anti-Cooper interaction, and no
reading that bears on either. The instrument measures a model electron's
coupling; it does not pair electrons, unbind pairs, or touch the machinery
by which gravity and the quantum are to be married. And the one place the
series' own apparatus even brushes the word Cooper cuts \emph{for} this
fence, not against it: the instrument's workshop keeps a small quantity of
retired, non-load-bearing material in what its record calls an attic ---
code that no reading passes through --- and the single fragment in the
entire apparatus that is allowed to compute without ever halting, the one
deliberately non-terminating specimen the finite gauge quarantines away
from everything that carries a reading, lives in exactly that attic, under
exactly that name. The device's only Cooper-adjacent artifact is the one
thing it has walled off from its own honest machinery on purpose. So
whatever the anti-Cooper-to-quantum-gravity idea may be worth as a
conjecture in someone's hands, it is worth nothing as a claim in this
book's, and the book makes none: it is an image at most, fenced hardest of
the three, and touching the device not at all. If a reader wants a road to
quantum gravity, this series does not sell one.

\section*{What the three houses share}

Take the chapter's discipline once more, because it is the volume's honesty
in miniature. Three problems larger than this instrument were named ---
the boundedness of fluids, the faintness of a particle the series does not
model, the marriage of gravity and the quantum --- and each was let
resonate exactly as far as a shared word or a shared shape allows, and
refused exactly at the point where a resonance would become a claim. The
instrument crossed no fence to reach any of them, and carries no theorem
away from any of them. A number can be remarkable and honest at once only
if it declines the larger stories it seems to invite; this chapter is where
this book declines them, by name, one house at a time. What the series does
owe --- the fluid face, and the deeper earnings that live in the device
rather than in prose --- the next chapter sets down plainly, alongside what
it has in fact read.
